% Texto original de esta revisión. Una sola leyenda para seis enunciados.
% La clave de lectura y el registro de relectura están en anexos/comprobaciones.tex.
\newcommand{\KalaContextoTexto}{%
  Documento M01: inspección ficticia de dos sensores. Los referentes se mantienen
  en todas sus páginas: 1 es el sensor A; 2 es el sensor B, distinto de A;
  3 es el registro de inspección; 4 es el evento de inspección, intervalo
  \([10,12]\); 5 es el evento de archivo, intervalo \([13,14]\).
  El tiempo se mide en segundos desde un mismo origen del experimento.
  C3 contiene exactamente esos dos sensores. P1 significa rojo según la
  clasificación declarada por el protocolo; P2 significa activo en la
  inspección. Ambas propiedades se evalúan en \(t_0=11\), bajo las mismas
  condiciones de observación. Las afirmaciones atribuidas al registro no se
  convierten por sí solas en hechos certificados.
}
\newcommand{\KalaEnunciadoUno}{\KalaFuente{\KalaReferencia{3}}{\KalaPosee{\KalaReferencia{1}}{2}}}
\newcommand{\KalaEnunciadoDos}{\KalaFuente{\KalaReferencia{3}}{\KalaNegacion{\KalaPosee{\KalaReferencia{2}}{2}}}}
\newcommand{\KalaEnunciadoTres}{\KalaFuente{\KalaReferencia{3}}{\KalaConjuncion{\KalaPosee{\KalaReferencia{1}}{1}}{\KalaNegacion{\KalaPosee{\KalaReferencia{2}}{1}}}}}
\newcommand{\KalaEnunciadoCuatro}{\KalaDefinicion{F}{\KalaSeleccion{x}{\KalaConcepto{3}}{\KalaConjuncion{\KalaPosee{\KalaVariable{x}}{1}}{\KalaPosee{\KalaVariable{x}}{2}}}}}
\newcommand{\KalaEnunciadoCinco}{\KalaPregunta{\KalaExiste{x}{\KalaMat{F}}{\KalaMiembro{\KalaVariable{x}}{\KalaConcepto{3}}}}}
\newcommand{\KalaEnunciadoSeis}{\KalaRelacion{1}{\KalaReferencia{4}\KalaSeparador\KalaReferencia{5}}}
\newcommand{\KalaMostrarEnunciado}[2]{%
  \par\noindent{\small\sffamily\color{KalaSuave}#1}\par
  \KalaMostrar{#2}\par\vspace{2mm}%
}
